% ---
% title: Variational Method
% date: 2026/1/12
% tag: ML
% ---
\section{變分法}

\textbf{变分法} （或称\textbf{变分法} ）是\href{https://en.wikipedia.org/wiki/Mathematical_analysis}{数学分析}的一个领域，利用变分，变分是\href{https://en.wikipedia.org/wiki/Function_(mathematics}{函数})中的微小变化 以及\href{https://en.wikipedia.org/wiki/Functional_(mathematics}{泛函})，用于求函数的极大值和最小值：从\href{https://en.wikipedia.org/wiki/Function_(mathematics}{一组函数})映射到\href{https://en.wikipedia.org/wiki/Real_number}{实数}的\href{https://en.wikipedia.org/wiki/Map_(mathematics}{映射})。\href{https://en.wikipedia.org/wiki/Calculus_of_variations#cite_note-2}{[a]} 泛函通常表示为包含函数及其\href{https://en.wikipedia.org/wiki/Derivative}{导}数的\href{https://en.wikipedia.org/wiki/Definite_integral}{确定积分}。最大化或最小泛函的函数可以通过变分法的\href{https://en.wikipedia.org/wiki/Euler–Lagrange_equation}{欧拉-拉格朗日方程}找到。

\subsection{例子}

一个简单的例子是找到连接两点的==最短长度曲线==。如果没有约束，解就是点之间的\href{https://en.wikipedia.org/wiki/Straight_line}{直线}。然而，如果曲线被限制在空间中的某个曲面上，那么解就不那么明显，可能存在许多解。这种解被称为\emph{\href{https://en.wikipedia.org/wiki/Geodesic}{测地线}}。\href{https://en.wikipedia.org/wiki/Fermat's_principle}{费马原理}提出了一个相关问题：光沿着连接两点的最短\href{https://en.wikipedia.org/wiki/Optical_length}{光学长度}路径前行，这取决于介质的材料。\href{https://en.wikipedia.org/wiki/Mechanics}{力学}中对应的一个概念是\href{https://en.wikipedia.org/wiki/Principle_of_least_action}{最小作用量/静止作用量原理}。

\section{欧拉-泊松方程（Euler–Poisson equation）}

如果\(S\) 取决于的 \(y(x)\)的更高阶导数，即如果对于

\[
S=\int_a^b{f(x,y(x),y'(x),\cdots y^{(n)}(x))}dx
\]

则 \(y\) 必须满足欧拉-\href{https://en.wikipedia.org/wiki/Siméon_Denis_Poisson}{泊松}方程，

\[
\frac{\partial f}{\partial y}-\frac{\partial}{\partial x}\frac{\partial f}{\partial y'}+\frac{\partial^2}{\partial x^2}\frac{\partial f}{\partial y''}-\cdots+(-1)^n\cdot \frac{\partial^n}{\partial x^n}\left[\frac{\partial f}{\partial y^{(n)}}\right]=0
\]

\subsection{示例}

求连接两点\(( x_1 , y_1)\) 和\(( x_2 , y_2)\) 的最短曲线\(y\)的表达式, 曲线的\href{https://en.wikipedia.org/wiki/Arc_length}{弧长}为:

\[
A[y]=\int_{x_1}^{x_2}{\sqrt{1+[y'(x)]^2}}dx
\]

所以

\[
\frac{\partial f}{\partial y}=\frac{\partial}{\partial x}\frac{\partial f}{\partial y'}=0
\]

即

\[
\frac{\partial}{\partial x}\left[\frac{y'(x)}{\sqrt{1+[y'(x)]^2}}\right]=0
\]

\[
\frac{y'(x)}{\sqrt{1+[y'(x)]^2}}=C \\
[y'(x)]^2=C(1+[y'(x)]^2) \\
\]

即有

\[
y'(x)=C
\]

两点之间直线最短，特别地，在本例中使用的一阶情况为\textbf{欧拉-拉格朗日方程}.
